\documentclass[11pt, a4paper]{article}

% bfw-style.tex — gemeinsame Style-Definitionen für alle Handouts/Aufgabenhefte
% Voraussetzung: Kompilation mit xelatex (wegen fontspec)

\usepackage[ngerman]{babel}
% Babel-ngerman macht " zu einem aktiven Shorthand (z.B. für "a -> ä). In unseren
% Texten verwenden wir " als normales Zeichen — daher Shorthand komplett ausschalten.
\shorthandoff{"}
\usepackage{geometry}
\geometry{a4paper, top=2.4cm, bottom=2.4cm, left=2.3cm, right=2.3cm,
          headheight=15pt, headsep=0.7cm, footskip=1.1cm}

\usepackage{fontspec}
% Fallback-Kette: Source Sans Pro -> Calibri -> Segoe UI -> Latin Modern Sans
% (Latin Modern ist immer mit MiKTeX vorhanden)
\IfFontExistsTF{Source Sans Pro}{%
  \setmainfont{Source Sans Pro}[Ligatures=TeX]%
}{\IfFontExistsTF{Calibri}{%
  \setmainfont{Calibri}[Ligatures=TeX]%
}{\IfFontExistsTF{Segoe UI}{%
  \setmainfont{Segoe UI}[Ligatures=TeX]%
}{%
  \setmainfont{Latin Modern Sans}[Ligatures=TeX]%
}}}
\IfFontExistsTF{Source Code Pro}{%
  \setmonofont{Source Code Pro}[Scale=0.92]%
}{\IfFontExistsTF{Consolas}{%
  \setmonofont{Consolas}[Scale=0.92]%
}{%
  \setmonofont{Latin Modern Mono}[Scale=0.92]%
}}

% Gesperrte Versalien für Labels/Titelbalken (Kapitälchen-Ersatz, funktioniert
% mit jeder OpenType-Schrift — Latin Modern Sans hat keine echten Kapitälchen)
\newcommand{\bfwlabelfont}{\footnotesize\bfseries\addfontfeatures{LetterSpace=8.0}}

\usepackage{xcolor}
% Farbpalette: hell, freundlich, modern
\definecolor{bfwPrimary}{HTML}{0EA5A4}    % Petrol/Türkis - Primär
\definecolor{bfwPrimaryDark}{HTML}{0F766E} % dunkler Petrol für Headings
\definecolor{bfwAccent}{HTML}{F59E0B}     % Amber - Akzent
\definecolor{bfwSuccess}{HTML}{10B981}    % Grün - Erfolg / Lösung
\definecolor{bfwWarn}{HTML}{F97316}       % Orange - Achtung
\definecolor{bfwInk}{HTML}{0F172A}        % fast Schwarz
\definecolor{bfwInkSoft}{HTML}{475569}    % gedämpftes Grau
\definecolor{bfwBgSoft}{HTML}{F8FAFC}     % off-white
\definecolor{bfwBgTeal}{HTML}{ECFEFF}     % helles Türkis
\definecolor{bfwBgAmber}{HTML}{FEF3C7}    % helles Gelb
\definecolor{bfwBgRose}{HTML}{FFE4E6}     % helles Rosé für Eselsbrücken
\definecolor{bfwTabHead}{HTML}{E4F3F2}    % zarte Kopfzeilen-Hinterlegung für Tabellen

\usepackage[colorlinks=true, linkcolor=bfwPrimaryDark, urlcolor=bfwPrimaryDark]{hyperref}
\usepackage{lastpage}
\usepackage{enumitem}
\setlist[itemize]{leftmargin=*, itemsep=2pt, topsep=2pt}
\setlist[enumerate]{leftmargin=*, itemsep=2pt, topsep=2pt}

\usepackage[most]{tcolorbox}
\usepackage{booktabs}
\usepackage{colortbl}
\usepackage{tikz}
\usetikzlibrary{decorations.pathmorphing, positioning, shapes.geometric, arrows.meta}

% ---------------------------------------------------------------------------
% Einheitliches Tabellen-Design
%   - etwas mehr Zeilenluft, dezente graue Linien (wirkt auch mit \hline ruhiger)
%   - Kopfzeile einer Tabelle bei Bedarf zart hinterlegen: \rowcolor{bfwTabHead}
% ---------------------------------------------------------------------------
\renewcommand{\arraystretch}{1.25}
\arrayrulecolor{bfwInkSoft!50}
\setlength{\heavyrulewidth}{1pt}
\setlength{\lightrulewidth}{0.5pt}

% ---------------------------------------------------------------------------
% Boxen — klare farbige Titelbalken mit gesperrten Versal-Labels
% (Titel bleibt per Option überschreibbar: \begin{merksatz}[title={...}])
% ---------------------------------------------------------------------------
\tcbset{bfwbox/.style={%
  enhanced, breakable,
  boxrule=0.5pt, arc=2.5pt,
  left=10pt, right=10pt, top=8pt, bottom=8pt,
  toptitle=4pt, bottomtitle=4pt,
  titlerule=0pt,
  fonttitle=\bfwlabelfont,
  coltitle=white
}}

% Merkkästchen — wichtige Definition / Kernpunkt
\newtcolorbox{merksatz}[1][]{%
  bfwbox,
  colback=bfwBgTeal,
  colframe=bfwPrimaryDark,
  colbacktitle=bfwPrimaryDark,
  title={MERKE},
  #1
}

% Eselsbrücke — Mnemoniks
\newtcolorbox{eselsbruecke}[1][]{%
  bfwbox,
  colback=bfwBgRose,
  colframe=bfwAccent!85!black,
  colbacktitle=bfwAccent!85!black,
  title={ESELSBRÜCKE},
  #1
}

% Beispiel-Box
\newtcolorbox{beispiel}[1][]{%
  bfwbox,
  colback=bfwBgSoft,
  colframe=bfwInkSoft,
  colbacktitle=bfwInkSoft,
  title={BEISPIEL},
  #1
}

% Lösungs-Box (für Aufgabenhefte)
\newtcolorbox{loesung}[1][]{%
  bfwbox,
  colback=bfwSuccess!7,
  colframe=bfwSuccess!65!black,
  colbacktitle=bfwSuccess!65!black,
  title={LÖSUNG},
  #1
}

% Achtung-Box — typische Fehler / Stolperfallen
\newtcolorbox{achtung}[1][]{%
  bfwbox,
  colback=bfwWarn!6,
  colframe=bfwWarn!85!black,
  colbacktitle=bfwWarn!85!black,
  title={ACHTUNG},
  #1
}

% Formel-Box — für zentrale Formeln (groß, ruhig, ohne Titelbalken)
\newtcolorbox{formelbox}[1][]{%
  enhanced,
  colback=bfwBgTeal,
  colframe=bfwPrimaryDark,
  boxrule=1pt, arc=3pt,
  left=10pt, right=10pt, top=6pt, bottom=6pt,
  #1
}

% Glossar-Eintrag
\newcommand{\glossareintrag}[2]{%
  \par\noindent\textbf{\color{bfwPrimaryDark}#1} \enspace #2\par\smallskip
}

% ---------------------------------------------------------------------------
% Abschnitts-Design: farbiges Nummern-Quadrat + feine Linie unter \section,
% dezent farbige \subsection/\subsubsection
% ---------------------------------------------------------------------------
\usepackage{titlesec}
\newcommand{\bfwSecBadge}[1]{%
  \begin{tikzpicture}[baseline={([yshift=-0.62em]n.center)}]
    \node[fill=bfwPrimaryDark, text=white, font=\normalsize\bfseries,
          minimum width=1.55em, minimum height=1.55em, inner sep=1pt,
          rounded corners=1.5pt] (n) {#1};
  \end{tikzpicture}}
\titleformat{\section}
  {\Large\bfseries\color{bfwPrimaryDark}}
  {\bfwSecBadge{\thesection}}{0.6em}{}
  [\vspace{-0.2em}{\color{bfwPrimary!60}\titlerule[0.8pt]}]
\titleformat{\subsection}
  {\large\bfseries\color{bfwPrimaryDark}}
  {\textcolor{bfwPrimary}{\thesubsection}}{0.5em}{}
\titleformat{\subsubsection}
  {\normalsize\bfseries\color{bfwInk}}
  {\textcolor{bfwPrimary}{\thesubsubsection}}{0.4em}{}
\titlespacing*{\section}{0pt}{1.6em}{1em}
\titlespacing*{\subsection}{0pt}{1.2em}{0.5em}

% TikZ-Stil "skizzenhaft" — etwas weicher als Rough.js, aber stabil
\tikzset{
  roughbox/.style = {
    draw=bfwPrimaryDark, thick, fill=bfwBgTeal,
    rounded corners=4pt, inner sep=6pt
  },
  roughaccent/.style = {
    draw=bfwAccent, thick, fill=bfwBgAmber,
    rounded corners=4pt, inner sep=6pt
  },
  roughline/.style = {
    draw=bfwInkSoft, thick, -{Stealth[length=2.5mm]}
  }
}

% Bullet-Symbole (bewusst kein FontAwesome — vermeidet Font-Installations-Probleme)
\newcommand{\faLightbulb}{$\bullet$}
\newcommand{\faBookmark}{$\bullet$}

% ---------------------------------------------------------------------------
% Kopf-/Fußzeile
%   Kopf links:  Wortlogo „Wirtschaftsfachwirt"
%   Kopf rechts: Dokument-Kurztext, gesetzt via \kopfzeile{...}
%   Fuß links:   © Can Siebert · 2026 — Fuß rechts: Seite X von Y
%   Titelseite (Seite 1) bleibt ohne Kopf-/Fußzeile (\handouttitel setzt
%   \thispagestyle{empty}).
% ---------------------------------------------------------------------------
\usepackage{fancyhdr}
\newcommand{\bfwKopfzeileText}{}
\newcommand{\kopfzeile}[1]{\renewcommand{\bfwKopfzeileText}{#1}}
\pagestyle{fancy}
\fancyhf{}
\fancyhead[L]{\small\bfseries\color{bfwPrimaryDark}Wirtschaftsfachwirt}
\fancyhead[R]{\small\color{bfwInkSoft}\bfwKopfzeileText}
\renewcommand{\headrulewidth}{0.6pt}
\renewcommand{\headrule}{{\color{bfwPrimary}\hrule width\headwidth height 0.6pt}}
\fancyfoot[L]{\small\color{bfwInkSoft}\copyright\ Can Siebert\,\textperiodcentered\,2026}
\fancyfoot[R]{\small\color{bfwInkSoft}Seite \thepage\ von \pageref*{LastPage}}
% Auch \thispagestyle{plain} (z.B. durch \maketitle) soll leer bleiben:
\fancypagestyle{plain}{\fancyhf{}\renewcommand{\headrulewidth}{0pt}\renewcommand{\headrule}{}}

% ---------------------------------------------------------------------------
% \handouttitel{Obertitel}{Untertitel}{Kurs-Label}{Termin-Zeile}
%   Einheitlicher professioneller Titelblock: Dachzeile, farbige Headline,
%   Untertitel, farbige Trennlinie und dezente Meta-Zeile (Kurs/Termin/Dozent).
%   Ersetzt \title/\maketitle. Direkt nach \begin{document} aufrufen.
% ---------------------------------------------------------------------------
\newcommand{\handouttitel}[4]{%
  \thispagestyle{empty}%
  \begingroup
  \setlength{\parindent}{0pt}%
  {\bfwlabelfont\color{bfwPrimary}WIRTSCHAFTSFACHWIRT (IHK)\par}
  \vspace{0.55em}
  {\huge\bfseries\color{bfwPrimaryDark}#1\par}
  \vspace{0.5em}
  {\large\color{bfwInkSoft}#2\par}
  \vspace{0.9em}
  {\color{bfwPrimary}\rule{\linewidth}{2pt}}\par
  \vspace{0.55em}
  {\small\color{bfwInkSoft}#3\ \,\textperiodcentered\,\ #4\ \,\textperiodcentered\,\ Dozent: Can Siebert\par}
  \endgroup
  \vspace{1.4em}%
}


\title{\vspace*{-1.5cm}\Huge\bfseries\color{bfwPrimaryDark}Aufgabenheft Lektion~4 (BWL)\\[0.4em]
  \large\color{bfwInkSoft}\textnormal{Unternehmenszusammenschlüsse --- Übungen im IHK-Klausurformat}}
\author{\textsf{Wirtschaftsfachwirt --- Betriebswirtschaftliche Grundlagen \textperiodcentered{} \small\copyright\ Can Siebert\,\textperiodcentered\,2027}}
\date{Selbstlernmaterial}

\begin{document}
\maketitle
\thispagestyle{empty}

\begin{merksatz}[title={\bfseries So nutzen Sie dieses Aufgabenheft}]
Bearbeiten Sie alle Aufgaben zunächst \emph{ohne} in die Lösungen zu schauen. Schreiben
Sie Ihre Antworten in vollständigen Sätzen --- die IHK-Prüfung verlangt formulierte
Begründungen, nicht bloße Stichworte. Beachten Sie die \emph{Operatoren} (nennen,
erläutern, beurteilen, begründen) --- sie geben den geforderten Antwort-Tiefegrad vor.
Erst danach öffnen Sie den Lösungsteil ab Seite~\pageref{loesungen} und vergleichen.
Ergänzend: Übungsband Betriebswirtschaft, Übungen 54--58 (mit Lösungen im Auszug).
\end{merksatz}

\tableofcontents
\vspace{1em}
\hrule

\section{Aufgaben}

\subsection*{Aufgabe~1 --- Gründe und Ziele von Zusammenschlüssen (8 Punkte)}

Die Geschäftsführung eines mittelständischen Maschinenbauers prüft einen Zusammenschluss
mit einem Wettbewerber, da „die Entwicklungskosten allein nicht mehr zu stemmen sind``.

\begin{enumerate}[label=(\alph*)]
  \item \textbf{Erklären} Sie, warum Unternehmen Zusammenschlüsse eingehen, und
    \textbf{unterscheiden} Sie dabei internes und externes Wachstum. \hfill (2~P.)
  \item \textbf{Nennen und erläutern} Sie vier typische Ziele von
    Unternehmenszusammenschlüssen. \hfill (6~P.)
\end{enumerate}

\subsection*{Aufgabe~2 --- Kooperation vs. Konzentration (10 Punkte)}

\begin{enumerate}[label=(\alph*)]
  \item \textbf{Unterscheiden} Sie Kooperation und Konzentration anhand der rechtlichen
    und der wirtschaftlichen Selbstständigkeit der beteiligten Unternehmen. \hfill (4~P.)
  \item \textbf{Ordnen} Sie die folgenden Formen der Kooperation bzw. Konzentration zu:
    Konsortium, Fusion, strategische Allianz, Konzern, Interessengemeinschaft, Trust.
    \hfill (3~P.)
  \item \textbf{Begründen} Sie, warum der Konzern zur Konzentration zählt, obwohl die
    Tochtergesellschaften rechtlich selbstständig bleiben. \hfill (3~P.)
\end{enumerate}

\subsection*{Aufgabe~3 --- Formen der Kooperation (12 Punkte)}

\textbf{Beschreiben} Sie die folgenden Kooperationsformen jeweils mit ihren Merkmalen
(Dauer, Zweck, Selbstständigkeit) und einem typischen Praxisbeispiel:

\begin{enumerate}[label=(\alph*)]
  \item Arbeitsgemeinschaft bzw. Konsortium \hfill (3~P.)
  \item Interessengemeinschaft \hfill (3~P.)
  \item Joint Venture (Gemeinschaftsunternehmen) \hfill (3~P.)
  \item Strategische Allianz \hfill (3~P.)
\end{enumerate}

\subsection*{Aufgabe~4 --- Zuordnung: horizontal, vertikal, diagonal (10 Punkte)}

\textbf{Ordnen} Sie die folgenden Zusammenschlüsse der zutreffenden Richtung
(horizontal, vertikal vor-/nachgelagert, diagonal) zu und \textbf{begründen} Sie kurz:

\begin{enumerate}[label=(\alph*)]
  \item Eine Brauerei übernimmt eine Mälzerei (Lieferant von Braumalz). \hfill (2~P.)
  \item Zwei regionale Molkereien schließen sich zusammen. \hfill (2~P.)
  \item Ein Automobilhersteller beteiligt sich mehrheitlich an einer Hotelkette. \hfill (2~P.)
  \item Ein Möbelhersteller übernimmt eine Einzelhandelskette für Möbel. \hfill (2~P.)
  \item Ein Großküchenhersteller kooperiert mit einer Geschäftsbank. \hfill (2~P.)
\end{enumerate}

\subsection*{Aufgabe~5 --- Konzern (12 Punkte)}

Die Alpha AG hält 75~\% der Anteile der Beta GmbH; zusätzlich besteht ein
Beherrschungsvertrag zwischen beiden Gesellschaften.

\begin{enumerate}[label=(\alph*)]
  \item \textbf{Definieren} Sie den Begriff Konzern. \hfill (3~P.)
  \item \textbf{Erläutern} Sie die Stellung von Mutter- und Tochtergesellschaft im
    Hinblick auf rechtliche und wirtschaftliche Selbstständigkeit. \hfill (4~P.)
  \item \textbf{Erklären} Sie Inhalt und Wirkung eines Beherrschungsvertrags. \hfill (3~P.)
  \item \textbf{Unterscheiden} Sie Unterordnungs- und Gleichordnungskonzern. \hfill (2~P.)
\end{enumerate}

\subsection*{Aufgabe~6 --- Fusion und Trust (8 Punkte)}

\begin{enumerate}[label=(\alph*)]
  \item \textbf{Erläutern} Sie die beiden Wege der Verschmelzung (durch Aufnahme, durch
    Neugründung) und \textbf{stellen} Sie beide schematisch \textbf{dar}. \hfill (4~P.)
  \item \textbf{Grenzen} Sie die Fusion vom Konzern \textbf{ab}. \hfill (2~P.)
  \item \textbf{Erklären} Sie den Begriff Trust. \hfill (2~P.)
\end{enumerate}

\subsection*{Aufgabe~7 --- Fallaufgabe: Kartell vs. Konzern (14 Punkte)}

\begin{beispiel}[title={Ausgangssituation}]
\textbf{Fall A:} Die drei größten Zementhersteller einer Region vereinbaren vertraglich
einheitliche Mindestpreise und teilen die Liefergebiete unter sich auf. Alle drei bleiben
rechtlich und wirtschaftlich selbstständig.

\textbf{Fall B:} Die Gamma AG erwirbt 100~\% der Anteile zweier Zementhersteller und
stellt beide unter ihre einheitliche Leitung. Beide Unternehmen bestehen als GmbH fort.
\end{beispiel}

\begin{enumerate}[label=(\alph*)]
  \item \textbf{Ordnen} Sie beide Fälle der zutreffenden Zusammenschlussform zu und
    \textbf{begründen} Sie Ihre Zuordnung anhand der Merkmale. \hfill (6~P.)
  \item \textbf{Beurteilen} Sie Fall~A kartellrechtlich: einschlägige Norm,
    Kartellart(en) und Rechtsfolge. \hfill (4~P.)
  \item \textbf{Erläutern} Sie, welcher staatlichen Kontrolle Fall~B unterliegen kann
    und unter welchen Voraussetzungen diese eingreift. \hfill (4~P.)
\end{enumerate}

\subsection*{Aufgabe~8 --- Fusionskontrolle (10 Punkte)}

\begin{enumerate}[label=(\alph*)]
  \item \textbf{Erklären} Sie Zweck und Grundprinzip der Fusionskontrolle nach dem GWB
    (Anmeldepflicht, Untersagungsgrund). \hfill (5~P.)
  \item \textbf{Nennen} Sie die zuständige deutsche Behörde und \textbf{erläutern} Sie,
    wann stattdessen die EU-Kommission zuständig ist. \hfill (3~P.)
  \item \textbf{Erklären} Sie den Begriff Ministererlaubnis. \hfill (2~P.)
\end{enumerate}

\subsection*{Aufgabe~9 --- Richtig oder falsch? (8 Punkte)}

\textbf{Beurteilen} Sie die folgenden Aussagen (richtig/falsch) mit kurzer Begründung:

\begin{enumerate}[label=(\alph*)]
  \item Bei der Kooperation geht die wirtschaftliche Selbstständigkeit der beteiligten
    Unternehmen verloren. \hfill (2~P.)
  \item Kartelle sind verboten, wenn sie zu spürbaren Wettbewerbsbeschränkungen führen.
    \hfill (2~P.)
  \item Bei einem Konzern werden rechtlich selbstständige Unternehmen unter einheitliche
    Leitung gestellt. \hfill (2~P.)
  \item Konsortien sind eine Form der Kooperation, die von der zuständigen
    Kartellbehörde grundsätzlich erlaubt wird. \hfill (2~P.)
\end{enumerate}

\newpage
\section{Lösungen}\label{loesungen}

\begin{loesung}[title={Lösung Aufgabe~1 --- Gründe und Ziele}]
\textbf{(a)} Unternehmen wollen wachsen und ihre Wettbewerbsposition verbessern.
\emph{Internes Wachstum} erfolgt aus eigener Kraft (Investitionen, eigene Innovationen),
\emph{externes Wachstum} durch Verbindung mit anderen Unternehmen --- ein Zusammenschluss
ermöglicht Ziele, die allein nicht oder nur langsam erreichbar wären (hier: Teilung der
Entwicklungskosten).

\textbf{(b)} Zum Beispiel: \emph{Rationalisierung/Kostensenkung} --- Größenvorteile,
gemeinsamer Einkauf mit Mengenrabatten, gemeinsame Verwaltung; \emph{Synergieeffekte} ---
Bündelung von Know-how, Teilung der F\&E-Kosten; \emph{Marktmacht} --- größere
Marktanteile, stärkere Position gegenüber Lieferanten und Abnehmern;
\emph{Risikostreuung} --- Verteilung des Risikos auf mehrere Geschäftsfelder oder
Projekte. (Ebenfalls richtig: Sicherung von Beschaffung/Absatz, Kapitalkraft,
Markterschließung, Schutz vor Übernahme.)
\end{loesung}

\begin{loesung}[title={Lösung Aufgabe~2 --- Kooperation vs. Konzentration}]
\textbf{(a)} \emph{Kooperation:} Die Unternehmen bleiben rechtlich \emph{und}
wirtschaftlich selbstständig; sie schränken ihre Entscheidungsfreiheit nur freiwillig im
vereinbarten Kooperationsbereich ein. \emph{Konzentration:} Die wirtschaftliche
Selbstständigkeit geht verloren (einheitliche Leitung); bei der Fusion und beim Trust
verliert mindestens ein Unternehmen zusätzlich die rechtliche Selbstständigkeit.

\textbf{(b)} Kooperation: Konsortium, strategische Allianz, Interessengemeinschaft.
Konzentration: Fusion, Konzern, Trust.

\textbf{(c)} Im Konzern unterstehen die Töchter der \emph{einheitlichen Leitung} der
Muttergesellschaft (Mehrheitsbeteiligung, ggf. Beherrschungsvertrag). Sie können ihre
unternehmerischen Entscheidungen nicht mehr frei treffen --- die \emph{wirtschaftliche}
Selbstständigkeit ist verloren. Entscheidend für die Einordnung als Konzentration ist der
Verlust der wirtschaftlichen, nicht der rechtlichen Selbstständigkeit.
\end{loesung}

\begin{loesung}[title={Lösung Aufgabe~3 --- Formen der Kooperation}]
\textbf{(a)} \emph{Arbeitsgemeinschaft/Konsortium:} befristeter Zusammenschluss
(i.\,d.\,R. als GbR) zur Durchführung \emph{eines} Auftrags oder Projekts; nach
Projektende Auflösung; volle Selbstständigkeit der Partner. Beispiele: ARGE mehrerer
Bauunternehmen für einen Tunnelbau; Bankenkonsortium zur Emission einer Anleihe.

\textbf{(b)} \emph{Interessengemeinschaft:} vertragliche, auf Dauer angelegte
Zusammenarbeit in einzelnen Teilbereichen (z.\,B. gemeinsame Forschung, Patentverwertung),
oft mit Gewinn- oder Kostenpool; am Markt treten die Unternehmen weiter getrennt auf.

\textbf{(c)} \emph{Joint Venture:} Die Partner gründen ein rechtlich selbstständiges
\emph{Gemeinschaftsunternehmen}, an dem sie Kapital, Führung und Risiko gemeinsam tragen
(oft 50:50). Beispiel: deutscher Hersteller gründet mit lokalem Partner eine gemeinsame
Produktionsgesellschaft zur Erschließung eines Auslandsmarkts.

\textbf{(d)} \emph{Strategische Allianz:} Zusammenarbeit --- häufig zwischen
Konkurrenten --- in einzelnen strategischen Geschäftsfeldern ohne gemeinsames
Tochterunternehmen; außerhalb der Allianz bleiben die Partner Wettbewerber. Beispiel:
Luftfahrt-Allianz mit gemeinsamem Streckennetz und Vielfliegerprogramm.
\end{loesung}

\begin{loesung}[title={Lösung Aufgabe~4 --- horizontal, vertikal, diagonal}]
\textbf{(a)} \emph{Vertikal (vorgelagert / Rückwärtsintegration):} Die Mälzerei ist
Lieferant der Brauerei --- vorgelagerte Wirtschaftsstufe; Ziel: Sicherung der
Rohstoffversorgung.

\textbf{(b)} \emph{Horizontal:} gleiche Branche, gleiche Produktionsstufe; Ziele:
Marktanteil, Größenvorteile, Rationalisierung.

\textbf{(c)} \emph{Diagonal (anorganisch/konglomerat):} branchenfremde Unternehmen ohne
Leistungszusammenhang; Ziel: Risikostreuung/Diversifikation.

\textbf{(d)} \emph{Vertikal (nachgelagert / Vorwärtsintegration):} Der Möbelhandel ist
die dem Hersteller nachgelagerte Absatzstufe; Ziel: Sicherung der Absatzwege.

\textbf{(e)} \emph{Diagonal:} Großküchenhersteller und Geschäftsbank gehören
verschiedenen Branchen ohne Produktionszusammenhang an (vgl. Übungsband, Übung 56).
\end{loesung}

\begin{loesung}[title={Lösung Aufgabe~5 --- Konzern}]
\textbf{(a)} Ein Konzern ist die Zusammenfassung eines herrschenden und eines oder
mehrerer abhängiger, \emph{rechtlich selbstständiger} Unternehmen unter
\emph{einheitlicher Leitung} (\S\,18 AktG).

\textbf{(b)} Beide Gesellschaften bleiben \emph{rechtlich} selbstständig (eigene Firma,
eigene Organe, eigene Bilanz). \emph{Wirtschaftlich} ist die Beta GmbH jedoch abhängig:
Die Alpha AG bestimmt über ihre Mehrheitsbeteiligung (75~\%) und den Beherrschungsvertrag
die Geschäftspolitik --- die wirtschaftliche Selbstständigkeit der Tochter ist verloren.

\textbf{(c)} Mit dem Beherrschungsvertrag (\S\,291 AktG) unterstellt sich eine
Gesellschaft der Leitung eines anderen Unternehmens. Die Mutter erhält ein direktes
\emph{Weisungsrecht} gegenüber der Geschäftsführung der Tochter; häufig wird er mit einem
Gewinnabführungsvertrag kombiniert. Im Gegenzug bestehen Schutzpflichten (z.\,B.
Verlustübernahme).

\textbf{(d)} \emph{Unterordnungskonzern:} herrschendes und abhängige(s) Unternehmen
(Mutter--Tochter, Regelfall). \emph{Gleichordnungskonzern:} rechtlich selbstständige
Unternehmen unter einheitlicher Leitung, ohne dass eines vom anderen abhängig ist.
\end{loesung}

\begin{loesung}[title={Lösung Aufgabe~6 --- Fusion und Trust}]
\textbf{(a)} \emph{Verschmelzung durch Aufnahme:} Unternehmen B überträgt sein gesamtes
Vermögen auf Unternehmen A und erlischt (A + B $\rightarrow$ A).
\emph{Verschmelzung durch Neugründung:} A und B übertragen ihr Vermögen auf eine neu
gegründete Gesellschaft C; beide erlöschen (A + B $\rightarrow$ C). Rechtsgrundlage ist
das Umwandlungsgesetz.

\textbf{(b)} Beim \emph{Konzern} bleiben alle Unternehmen rechtlich bestehen und nur die
wirtschaftliche Selbstständigkeit geht verloren; bei der \emph{Fusion} verliert
mindestens ein Unternehmen auch seine rechtliche Selbstständigkeit --- es entsteht ein
einziger Rechtsträger.

\textbf{(c)} Der \emph{Trust} bezeichnet die vollständige Verschmelzung ehemals
selbstständiger Unternehmen zu einer wirtschaftlichen und rechtlichen Einheit mit
marktbeherrschender Stellung (historisch v.\,a. USA; daher „Antitrust-Recht`` für das
Wettbewerbsrecht).
\end{loesung}

\begin{loesung}[title={Lösung Aufgabe~7 --- Fallaufgabe Kartell vs. Konzern}]
\textbf{(a)} \emph{Fall A: Kartell} --- vertragliche Absprache rechtlich und
wirtschaftlich selbstständig bleibender Unternehmen derselben Wirtschaftsstufe mit dem
Ziel der Wettbewerbsbeschränkung (Preise, Gebiete). \emph{Fall B: Konzern} --- die Gamma
AG stellt die beiden weiterhin rechtlich selbstständigen GmbHs über eine 100-\%-Beteiligung
unter ihre einheitliche Leitung; die wirtschaftliche Selbstständigkeit der Töchter geht
verloren.

\textbf{(b)} Die Vereinbarung verstößt gegen das \emph{Kartellverbot} des \S\,1 GWB
(bzw. Art.\,101 AEUV): Es liegt ein \emph{Preiskartell} (Mindestpreise) kombiniert mit
einem \emph{Gebietskartell} (Aufteilung der Liefergebiete) vor --- sogenannte
Hardcore-Beschränkungen, für die keine Freistellung in Betracht kommt. Rechtsfolge: Die
Vereinbarung ist nichtig; das Bundeskartellamt kann empfindliche Bußgelder verhängen.

\textbf{(c)} Fall B unterliegt der \emph{Fusionskontrolle} (\S\S\,35\,ff. GWB): Bei
Überschreiten der Umsatzschwellen muss der Zusammenschluss \emph{vor} dem Vollzug beim
Bundeskartellamt angemeldet werden. Untersagt wird er, wenn er wirksamen Wettbewerb
erheblich behindern würde, insbesondere durch Entstehung oder Verstärkung einer
marktbeherrschenden Stellung; möglich sind auch Freigaben unter Auflagen. Bei
gemeinschaftsweiter Bedeutung prüft stattdessen die EU-Kommission.
\end{loesung}

\begin{loesung}[title={Lösung Aufgabe~8 --- Fusionskontrolle}]
\textbf{(a)} \emph{Zweck:} Schutz des Wettbewerbs vor dauerhafter Ausschaltung durch
Unternehmenskonzentration. \emph{Grundprinzip:} präventive Kontrolle ---
Zusammenschlüsse oberhalb bestimmter Umsatzschwellen (u.\,a. weltweit zusammen mehr als
500~Mio.~€) müssen vor dem Vollzug angemeldet werden; der Vollzug ist bis zur Freigabe
verboten. \emph{Untersagungsgrund:} erhebliche Behinderung wirksamen Wettbewerbs,
insbesondere Entstehung oder Verstärkung einer marktbeherrschenden Stellung (Vermutung
ab 40~\% Marktanteil).

\textbf{(b)} Zuständig ist das \emph{Bundeskartellamt} (Sitz: Bonn). Bei
Zusammenschlüssen mit \emph{gemeinschaftsweiter Bedeutung} (sehr hohe, europaweite
Umsätze) ist nach der EU-Fusionskontrollverordnung ausschließlich die
\emph{EU-Kommission} zuständig („One-Stop-Shop``).

\textbf{(c)} Die \emph{Ministererlaubnis} (\S\,42 GWB) erlaubt ausnahmsweise einen vom
Bundeskartellamt untersagten Zusammenschluss, wenn gesamtwirtschaftliche Vorteile oder
ein überragendes Interesse der Allgemeinheit die Wettbewerbsbeschränkung überwiegen;
sie wird vom Bundeswirtschaftsministerium erteilt.
\end{loesung}

\begin{loesung}[title={Lösung Aufgabe~9 --- Richtig oder falsch?}]
\textbf{(a) Falsch.} Bei der Kooperation behalten die Unternehmen ihre wirtschaftliche
(und rechtliche) Selbstständigkeit; sie wird nur im Kooperationsbereich freiwillig
eingeschränkt.

\textbf{(b) Richtig.} Nach \S\,1 GWB sind wettbewerbsbeschränkende Vereinbarungen
verboten, wenn sie den Wettbewerb spürbar beschränken; nur eng begrenzte Ausnahmen
(z.\,B. bestimmte Mittelstandskooperationen) sind freigestellt.

\textbf{(c) Richtig.} Der Konzern fasst rechtlich selbstständige Unternehmen unter
einheitlicher Leitung zusammen (\S\,18 AktG); verloren geht die wirtschaftliche
Selbstständigkeit.

\textbf{(d) Richtig.} Konsortien (wie Arbeitsgemeinschaften) sind projektbezogene
Kooperationen, die den Wettbewerb regelmäßig nicht spürbar beschränken --- sie sind
kartellrechtlich grundsätzlich erlaubt (vgl. Übungsband, Übung 54).
\end{loesung}

\vspace{1em}
\hrule
\vspace{0.8em}
\noindent\textsf{\small Weiterüben: Übungsband Betriebswirtschaft, Übungen 54--58 (mit Lösungsteil) und das
interaktive Quiz dieser Lektion. Dies war die letzte Lektion --- viel Erfolg bei der IHK-Prüfung!}

\end{document}
